% ======================================
% 01-packages.tex - 宏包引入
% ======================================
% 注意：ctex 已在文档类中加载，这里不需要重复

% ========== 页面设置 ==========
\usepackage[a4paper, left=2.5cm, right=2.5cm, top=2.5cm, bottom=2.5cm]{geometry}
\usepackage{fancyhdr}

% ========== 数学宏包 ==========
\usepackage{amsmath, amssymb, amsthm, mathtools}
\usepackage{newpxmath}
\usepackage{esint}
\usepackage{bm}

% ========== 表格与排版 ==========
\usepackage{multirow}
\usepackage{diagbox}
\usepackage{needspace}
\usepackage{makecell}
\usepackage{booktabs}
\usepackage{etoolbox}
\usepackage{environ}  % 用于解答隐藏功能
\usepackage{enumitem} % 用于自定义列表

% ========== 盒子与高亮 ==========
\usepackage[most]{tcolorbox}
\usepackage{xcolor}
\usepackage{newunicodechar}

% ========== 图标支持 ==========
\usepackage{fontawesome5}  % 用于交互元素图标

% ========== 图表与 TikZ ==========
\usepackage{tikz}
\usetikzlibrary{arrows.meta}
\usetikzlibrary{patterns}
\usepackage{pgfplots}
\pgfplotsset{compat=1.18}
\usepgfplotslibrary{fillbetween}
\usepackage{tikz-3dplot}
\tdplotsetmaincoords{70}{110}
\usepackage{capt-of}
\usepackage{graphicx,caption,subcaption}
\renewcommand{\thesubfigure}{\Alph{subfigure}}

% ========== 超链接（可选） ==========
\@ifundefined{hypersetup}{
    \usepackage{hyperref}
}{}
